\chapter{Практика}

\section{Выбор данных для анализа}

Для анализа я выбрал небольшую базу данных с описанием флагов разных стран.

\section{Описание набора данных}
Набор данных содержит N строк, в каждой из которых записаны следующие данные:

\begin{enumerate}
\item name:	Name of the country concerned
\item landmass:	1=N.America, 2=S.America, 3=Europe, 4=Africa, 4=Asia, 6=Oceania
\item zone:	Geographic quadrant, based on Greenwich and the Equator; 1=NE, 2=SE, 3=SW, 4=NW area:	in thousands of square km
\item population:	in round millions
\item language: 1=English, 2=Spanish, 3=French, 4=German, 5=Slavic, 6=Other Indo-European, 7=Chinese, 8=Arabic, 9=Japanese/Turkish/Finnish/Magyar, 10=Others
\item religion: 0=Catholic, 1=Other Christian, 2=Muslim, 3=Buddhist, 4=Hindu, 5=Ethnic, 6=Marxist, 7=Others
\item bars:     Number of vertical bars in the flag
\item stripes:  Number of horizontal stripes in the flag
\item colours:  Number of different colours in the flag
\item red:      0 if red absent, 1 if red present in the flag
\item green:    same for green
\item blue:     same for blue
\item gold:     same for gold (also yellow)
\item white:    same for white
\item black:    same for black
\item orange:   same for orange (also brown)
\item mainhue:  predominant colour in the flag (tie-breaks decided by taking the topmost hue, if that fails then the most central hue, and if that fails the leftmost hue)
\item circles:  Number of circles in the flag
\item crosses:  Number of (upright) crosses
\item saltires: Number of diagonal crosses
\item quarters: Number of quartered sections
\item sunstars: Number of sun or star symbols
\item crescent: 1 if a crescent moon symbol present, else 0
\item triangle: 1 if any triangles present, 0 otherwise
\item icon:     1 if an inanimate image present (e.g., a boat), otherwise 0
\item animate:  1 if an animate image (e.g., an eagle, a tree, a human hand) present, 0 otherwise
\item text:     1 if any letters or writing on the flag (e.g., a motto or slogan), 0 otherwise
\item topleft:  colour in the top-left corner (moving right to decide tie-breaks)
\item botright: Colour in the bottom-left corner (moving left to decide tie-breaks)
\end{enumerate}

Пример:
\begin{lstlisting}
Afghanistan,5,1,648,16,10,2,0,3,5,
1,1,0,1,1,1,0,green,0,0,0,0,1,0,0,1,0,0,black,green
\end{lstlisting}

\section{Анализ данных}

\subsection{Считывание данных}
С помощью библиотеки pandas считаем данные из файла:

\begin{lstlisting}
    data = pd.read_csv(data_file, sep=",")
\end{lstlisting}

\subsection{Подсчет статистических метрик}
Библиотека pandas предоставляет удобный способ группировки данных и подсчета статистических метрик.
Сгруппируем страны по регионам и для каждого региона составим статистику о внешнем виде флагов (цвета, фигуры, изображения и т.д)
\begin{lstlisting}
info = data.groupby(["landmass"]).agg(
    #colors
    green=("green", "sum"),
    red=("red", "sum"),
    gold=("gold", "sum"),
    blue=("blue", "sum"),
    white=("white", "sum"),
    black=("black", "sum"),
    orange=("orange", "sum"),
    
    #figure
    circles=("circles", "sum"),
    sunstars=("sunstars", "sum"),
    quarters=("quarters", "sum"),
    crosses=("crosses", "sum"),
    saltires=("saltires", "sum"),
    animate=("animate", "sum"),
    icon=("icon", "sum"),
    
    #text
    text=("text", "sum"),
    
    #count
    count=("country", "count"),
     
)
\end{lstlisting}

Переменная info будет принадлежать типу данных DataFrame, который представляет из себя таблицу.
Можно вывести переменную info с помощью print(), а лучше записать в файл. Получим вот такую таблицу:

\begin{figure}[H]
    \centering
    \includegraphics{table.png}
    \caption{Таблица со статистикой}
\end{figure}

\section{Построение графиков}
Анализировать такие большие таблицы неудобно, поэтому с помощью библиотеки matplotlib построим диаграммы. \\
Отобразим распределнеие государств по регионам:

\begin{figure}[H]
    \centering
    \includegraphics{d1.png}
    \caption{Диаграмма распределения стран по регионам}
\end{figure}

Можно увидеть, что в Европе, Азии и Африки сконцентрированно наибольшее кол-во стран. \\

Для каждого региона построим по 2 диаграммы, которые будут показывать соотношение цветов флагов в регионе и соотношение фигур на флаге (крестки, кружки, звезды и т.д):

\begin{figure}[H]
    \centering
    \includegraphics{d2.png}
    \caption{Диаграммы соотношения цветов и фигур в регионах}
\end{figure}

\endinput